%Grecias, PCSI, pg 764
\begin{exercise}[Bouteille brisée]
Lorsque l'on place une bouteille de verre, pleine à \SI{95}{\percent} d'eau dehors en plein hiver, contenant \SI{1}{\liter} d'eau, il arrive qu'on la retrouve brisée le lendemain.
\begin{questions}
\question Avant de se briser, la bouteille contient une fraction $x$ de glace. Calculer le volume de glace et le volume d'eau dans la bouteille avant qu'elle ne se brise.
\question Quelle discontinuité de changement d'état peut-on déduire de cette observation ? Où doit-on en chercher une explication ?
\question Quel est le volume de glace une fois la bouteille brisée ?
\end{questions}
\textbf{Données à \SI{0}{\celsius} : } $V_m^\star (\ce{H2O(s)}) = 19.7 ~ \si{\centi\meter\cubed\per\mole}$, $V_m^\star (\ce{H2O(l)}) = 18.1 ~ \si{\centi\meter\cubed\per\mole}$

\end{exercise}